\section{开场与阅读范围}
\presenter{第8队}

\begin{frame}[plain]
    \vspace{4mm}{\small\color{accent}天津大学 NLP\quad 第 8 队}\par
    \vspace{8mm}{\fontsize{27}{34}\selectfont\bfseries 长上下文场景\\结构化记忆建模}\par
    \vspace{5mm}{\Large 六篇论文阅读与系统性比较}\par
    \vspace{6mm}\textcolor{accent}{\rule{\textwidth}{0.8pt}}\par
    \vspace{3mm}{\small 组织、访问、更新、利用与操作验证}
    \note{计划15秒。大家好，我们是第八队。这次阅读围绕长上下文场景中的结构化记忆，介绍六篇论文，并用统一的问题比较它们。我们关注历史怎样被保存、找到、更新并用于后续任务。}
\end{frame}

% 来源：阅读范围：六篇原论文；下列问题为综述整理
\begin{frame}{长历史中的关系、范围与状态}
    \large
    长期记忆需要保留关系、定位证据，并跟随新信息更新。\par
    \vspace{5mm}\normalsize
    \begin{tabularx}{\linewidth}{@{}p{26mm}Y@{}}
        \toprule
        环节 & 需要回答的问题 \\
        \midrule
        组织 & 人物、事件、时间和来源关系是否保留下来？\\[2mm]
        访问 & 如何定位正确的人物、事件与版本？\\[2mm]
        更新与利用 & 新旧状态如何区分，已取回的证据如何解释？\\[2mm]
        操作验证 & 改了哪条记忆，执行后的状态是否符合意图？\\
        \bottomrule
    \end{tabularx}
    \note{计划60秒。假设用户先说周五去A地，后来改成周六去B地。完整历史虽然包含两条信息，模型仍可能混淆哪个计划有效；如果只检索语义最相似的片段，也可能只找到旧计划。即使都找到了，回答时还可能把两个版本拼在一起。最后，如果用户要求暂时不要再提这次行程，系统还要判断是降低优先级还是删除。这个小例子把问题分成四层：组织是否保留关系，访问是否选对范围，更新和回答是否处理好证据，操作结果是否符合要求。后面的六篇工作分别围绕这些环节展开。}
\end{frame}

% 来源：时间为讲解计划；实际时长以本人排练与录制为准
\begin{frame}{汇报安排}
    \begin{tabularx}{\linewidth}{@{}p{25mm}Y r@{}}
        \toprule
        环节 & 内容 & 计划时间 \\
        \midrule
        开场 & 共同问题与阅读范围 & 2 分钟 \\
        六篇论文 & StructMem、Hindsight、Semantic XPath\newline HyMEM、MetaMem、Text2Mem & 每人约8分钟 \\
        系统性总结 & 机制比较、共同案例、实验结果与课设选择 & 10 分钟 \\
        \bottomrule
    \end{tabularx}\par
    \vspace{5mm}\key{2＋6×8＋10＝60 分钟}

    \note{计划45秒。先由六位成员分别说明问题、方法、证据和局限，每部分以八分钟为目标，再用十分钟回答这些工作之间是什么关系。StructMem和Hindsight主要讨论长对话记忆，Semantic XPath研究结构访问，HyMEM面向网页操作，MetaMem关注证据利用，Text2Mem关注操作语义与状态效果。最后结合各自的任务和实验条件，比较这些机制怎样发挥作用。下面进入StructMem。}
\end{frame}
